Rotation averaging yields
\((-\Delta_{T,\tau}+V(|\tau|))H(T,|\tau|)\ge0\).
We first bound the radial comparison solutions. For \(S=1\) or \(3\)
and fixed \(\kappa>0\), solve
\[
 u''+s^{-1}u'=(V+\kappa^2)u,\qquad b<s<S,
\]
with endpoint values \(0,1\) in either order.
Set \(x=\log(s/b)\), \(A=\log(S/b)\), and
\(P(x)=s^2(V(s)+\kappa^2)\). The equation in \(x\) is \(u_{xx}=Pu\).
Its solution \(\psi\) with \(\psi(0)=0\) and \(\psi'(0)=1\)
satisfies \(\psi(x)\ge x\): it begins positive, then has
\(\psi''\ge0\) wherever it is nonnegative. Integration gives
\[
 \frac{\psi(x)}x=1+\int_0^x \frac{z(x-z)}x
                  P(z)\frac{\psi(z)}z\,dz .
\]
For \(0\le z\le x\le A\) the coefficient \(z(x-z)/x\) is at most
\(\min(z,A-z)\). Gronwall's inequality therefore bounds
\(\psi(A)/A\) above by
\(\exp(\int_0^A\min(z,A-z)P(z)\,dz)\).
The part from \(V\) is at most
\(\int_b^S s\log(s/b)V(s)\,ds\le M\).
Writing \(t=A-z\), the part from \(\kappa^2\) is at most
\(\kappa^2S^2\int_0^\infty t e^{-2t}\,dt\).
Consequently the solution with values \(0,1\) is
\(\psi(x)/\psi(A)\ge c x/A\). Reflection \(x\mapsto A-x\)
preserves the weighted integral used in Gronwall's inequality
and yields \(u\ge c(1-x/A)\) for values \(1,0\).

Put \(m(s)=\int_0^3\sin(\pi T/3)H(T,s)\,dT\).
Integration by parts twice gives
\[
 \int_0^3\sin(\pi T/3)\,\partial_T^2H(T,s)\,dT
   =-(\pi/3)^2m(s)+(\pi/3)\big(H(0,s)+H(3,s)\big).
\]
The boundary term is nonnegative, so the supersolution inequality
implies \(m''+s^{-1}m'\le(V+\pi^2/9)m\).
Since \(m(3)\ge0\), the maximum principle and the radial
comparison with \(S=3\), \(\kappa=\pi/3\), and boundary values
\(1,0\) give \(m(s)\ge c(M)m(b)\log(3/s)/D\). In particular
\[
 \int_0^3\int_{3/4\le|\tau|\le5/4}H(T,|\tau|)\,d\tau\,dT
 \ge c(M)m(b)/D.
\]

We transfer this integrated bound to \(H(T,1)\) for \(|T|\le1\).
On \(-2<T<4,\ 1/2<|\tau|<3/2\) one has \(\Delta H\le MH\).
Write \(x=(T,\tau)\). The function
\(H(x)\cos(\sqrt{M+1}\,v)\) in the additional real variable \(v\)
is nonnegative and superharmonic on a fixed small slab about
\(v=0\). The mean inequality on the four-dimensional ball of
radius \(2\rho\) about \((x,0)\), restricted to
\(B_\rho(x)\times[-\rho,\rho]\), gives
\[
 H(x)\ge c\rho^{-3}\int_{B_\rho(x)} H(y)\,dy
\]
for one sufficiently small \(\rho>0\), uniformly on a neighborhood
of the compact connected band
\([-1,3]\times\{3/4\le|\tau|\le5/4\}\).
Let \(B_1,\ldots,B_N\) be a finite cover of the band by balls of
radius \(\rho/4\) centered in the band. The overlap graph is
connected because each ball meets the band and the band is
connected. The integrated bound gives one ball \(B_{j_0}\) with
mass \(I_{j_0}=\int_{B_{j_0}}H\ge c(M)m(b)/(ND)\).
For any overlapping pair \(B_i,B_j\), every \(x\in B_j\) has
\(B_i\subset B_\rho(x)\). Nonnegativity and the mean inequality
therefore give \(\inf_{B_j}H\ge c\rho^{-3}I_i\). Integrating on
\(B_j\) gives \(I_j\ge c\rho^{-3}|B_j|I_i\), with a fixed
positive factor because \(|B_j|\) is proportional to \(\rho^3\).
Thus along a chain from \(B_{j_0}\) the lower mass bound passes
from each known ball to the next one. A shortest chain has at
most \(N-1\) edges. Applying the mean inequality within each
ball to its resulting mass gives \(\inf_{B_j}H\ge c(M)m(b)/D\)
for every \(j\). In particular this holds at \((T,\tau)\) with
\(|T|\le1\) and \(|\tau|=1\).

For the return passage use the radial solution for
\(S=1\), \(\kappa=\pi/2\), \(u(b)=0\), \(u(1)=1\). On
\(|T|<1,\ b<|\tau|<1\), the function
\[
 w(T,\tau)=c(M)m(b)D^{-1}\cos(\pi T/2)\,u(|\tau|)
\]
solves \((-\Delta_{T,\tau}+V)w=0\).
It vanishes on \(|T|=1\) and \(|\tau|=b\); it is at most
\(H(T,1)\) on \(|\tau|=1\) for a sufficiently small \(c(M)\).
The maximum principle and the radial lower bound imply
\[
 H(0,r)\ge c(M)\frac{m(b)}D
              \frac{\log(r/b)}{\log(1/b)}.
\]
Use \(\log(1/b)\le D\) and the positive lower bound for
\(\sin(\pi T/3)\) on \([1,2]\) to conclude.
